% from ega4-17-fr.tex
\oldpage[IV-4]{105}le morphisme $V_1\to U_1$ restriction de $f_1$ soit un isomorphisme (17.15.14.1). Appliquant (8.9.1) et (8.10.5, (vi) et (x)), on voit donc d'abord qu'il existe une extension finie radicielle $k''$ de $k$ et un morphisme fini surjectif $Y''\to X''=X_{(k'')}$ tels que $Y_1=Y''\times_{X''}X_1=Y''\otimes_{k''}k_1$; puisque $Y_1$ est normal et $k_1$ parfait, il résulte de (6.7.7) que $Y''$ est géométriquement normal sur $k''$. Comme les projections $X_1\to X''$ et $Y_1\to Y''$ sont des morphismes entiers, surjectifs et radiciels, ce sont des homéomorphismes (2.4.5), et si $U''$ et $V''$ sont les images de $U_1$ et $V_1$ dans $X''$ et $Y''$ respectivement, ce sont des ouverts denses dans $X''$ et $Y''$ respectivement tels que $V''=f''^{-1}(U'')$, où $f'':Y''\to X''$ est le morphisme structural. Comme on a $U_1=U''\otimes_{k''}k_1$, il résulte de (8.10.5, (i)) qu'il existe une extension finie radicielle $k'$ de $k''$ telle que si l'on pose $X'=X\otimes_k k'=X''\otimes_{k''}k'$, $Y'=Y''\otimes_{k''}k'$, $f'=f''_{(k')}:Y'\to X'$, et si $U'$ et $V'$ sont les images de $U''$ et $V''$ dans $X'$ et $Y'$ respectivement, la restriction $V'\to U'$ de $f'$ soit un isomorphisme. Comme $Y'$ est normal et $f'$ entier et birationnel, on conclut de (17.15.14.1) que $Y'$ est isomorphe au normalisé de $X'$, ce qui prouve le lemme puisque $Y'$ est géométriquement normal.

Revenons maintenant à la démonstration de (17.15.14), et appliquons à $k$ et $X$ le lemme (17.15.14.2); puisque $\dim(X)\leq1$, on a aussi $\dim(X_{(k')})\leq1$ (4.1.4), et le normalisé $Y'$ de $X_{(k')}$ est aussi de dimension $\leq1$ (5.4.2 et II, 7.4.6). Dire que $Y'$ est géométriquement normal sur $k'$ revient alors à dire, en vertu des définitions et de (II, 7.4.5), que $Y'$ est géométriquement régulier sur $k'$, donc lisse sur $k'$ (17.5.1), ce qui prouve (17.15.14).

\begin{proposition}[17.15.15]
\label{IV.17.15.15}
Soit $f:X\to Y$ un morphisme localement de présentation finie. Pour que $f$ soit lisse en un point $x\in X$, il faut et il suffit que $f$ soit plat au point $x$ et que $\Omega^1_f$ soit un $\sh{O}_X$-Module localement libre dans un voisinage de $x$, de rang en $x$ égal à $\dim_x f$.
\end{proposition}

La nécessité des conditions résulte de (17.5.1) et (17.10.2). Inversement, si ces conditions sont vérifiées, et si $y=f(x)$, il suffit de montrer (17.5.1) que $f^{-1}(y)$ est lisse sur $k(y)$ au point $x$; mais cela résulte de la définition de $\dim_x f$ (17.10.1), de (16.4.5) et de (17.15.5).

\subsection{Quasi-sections de morphismes plats ou lisses.}
\EGAReaderTarget{ega4.c04.subsec.17-16}{subsection:IV.17.16}

Les énoncés de ce numéro complètent ceux de (14.5), moyennant des hypothèses de platitude.

\begin{proposition}[17.16.1]
\label{IV.17.16.1}
Soit $f:X\to S$ un morphisme plat localement de présentation finie. Soient $s\in S$, $x$ un point fermé de $X_s$ tel que $\sh{O}_{X_s,x}$ soit un anneau de Cohen--Macaulay; alors il existe un voisinage ouvert $U$ de $s$ dans $S$ et un sous-préschéma $X'\subset f^{-1}(U)$ de $X$ tel que $x\in X'$ et que le morphisme $X'\to U$, restriction de $f$, soit plat, quasi-fini et de présentation finie.
\end{proposition}

La question étant locale sur $X$ et $Y$, on peut supposer $f$ de présentation finie. Soit $(t_i)_{1\leq i\leq m}$ un système de paramètres de l'anneau local $\sh{O}_{X_s,x}$; il existe un voisinage ouvert affine $V=\Spec(A)$ de $x$ dans $X$ et $m$ sections $g_i$ de $\sh{O}_X$ au-dessus de $V$ telles que

